% =====================================================================
% セットアップ編 第1章 Gitを使えるようにする
% 元原稿: chapters/01_git.tex（本編から分離）
% =====================================================================
\chapter{Gitを使えるようにする}

この章では、\textbf{自分のPCでGitを使えるようにするまで}を扱います。やることは「Gitをインストールする」と「名前とメールアドレスを登録する」の2つだけです。設定はPC全体に保存されるので、\textbf{ここは一度やれば終わり}です。この章を終えたら、本編『はじめてのバージョン管理』に進んでください。


% ---------------------------------------------------------------------
\section{Gitのインストール}

まず、Gitが\textbf{すでに入っていないか}を確認しましょう。ターミナル（WindowsはPowerShellやコマンドプロンプト）で次を実行します。

\begin{codebox}[bash]
\begin{code}
git --version
# 出力例: git version 2.45.0
\end{code}
\end{codebox}

\textbf{この結果によって、次に進む先が変わります。}

\begin{itemize}[leftmargin=1.4em]
  \item \textbf{バージョンが表示された人}（例：\cmd{git version 2.45.0}）\\
        → Gitは\textbf{すでに入っています}。この後のOS別の手順は\textbf{すべて飛ばして}、\ref{sec:first-config}節\textbf{「最初の設定」}へ進んでください。
  \item \textbf{\cmd{command not found}（見つからない）と出た人}\\
        → Gitが\textbf{入っていません}。お使いのOSの手順へ進んでください。\\
        \textbf{Windows} → \ref{sec:install-win}項　／　\textbf{macOS} → \ref{sec:install-mac}項　／　\textbf{Linux} → \ref{sec:install-linux}項
\end{itemize}

インストールが終わったら、もう一度 \cmd{git -{}-version} を実行してください。バージョンが表示されれば成功です。

\begin{warnbox}
\textbf{インストール直後は、開いていたターミナルではまだ \cmd{git --version} が反映されないことがあります。} PATH（コマンドの置き場所を示す設定）がターミナルの起動時に読み込まれるためです。\cmd{command not found} のままなら、\textbf{ターミナルをいったん閉じて開き直して}から、もう一度実行してください。
\end{warnbox}

\subsection{Windows}
\label{sec:install-win}

Windowsでは「Git for Windows」をインストールします。これには \textbf{Git Bash} という、macOS/Linux風のコマンドが使えるターミナルが付いてきます。本シリーズのWindowsの説明は、この Git Bash を使う前提です。

\begin{enumerate}[leftmargin=1.6em]
  \item \url{https://git-scm.com/download/win} からインストーラをダウンロードし、ダブルクリックで起動します。
  \item インストーラの画面が続きます。\textbf{基本はすべてデフォルトのまま「Next」で進めて問題ありません。} ただし次の3つの画面だけは選択を確認してください。
\end{enumerate}

\noindent\textbf{① 既定のエディタ}

デフォルトはVimですが不慣れだと難しいため、VS Codeを使う予定なら \textbf{「Use Visual Studio Code as Git's default editor」} を選びます。
\screenshot[0.62]{01_git/win_install_05.png}{既定エディタの選択}

\noindent\textbf{② 初期ブランチ名}

\textbf{「Override the default branch name…」} を選び、\cmd{main} と入力します（世界的に \cmd{main} が標準）。
\screenshot[0.62]{01_git/win_install_06.png}{初期ブランチ名の設定}

\noindent\textbf{③ 改行コードの扱い}

WindowsとほかのOSで改行コードが異なるため重要です。複数環境で開発するなら \textbf{「Checkout as-is, commit Unix-style line endings」} を選びます。
\screenshot[0.62]{01_git/win_install_10.png}{改行コードの設定}

残りの画面はデフォルトのまま進め、最後に「Install」をクリックします。インストール後、スタートメニューから \textbf{Git Bash} を起動し、\cmd{git --version} でバージョンが出れば完了です。

\subsection{macOS}
\label{sec:install-mac}

次のどちらか一方でOKです。

\noindent\textbf{方法A: Xcode Command Line Tools（最も手軽）}

ターミナル（\cmd{アプリケーション > ユーティリティ > ターミナル.app}）を開き、次を実行します。

\begin{codebox}[bash]
\begin{code}
xcode-select --install
\end{code}
\end{codebox}

ダイアログが出たら「インストール」をクリックします。

\noindent\textbf{方法B: Homebrew（最新版を使いたい場合）}

Homebrew（\url{https://brew.sh/index_ja}）が入っていれば、次で最新のGitが入ります。

\begin{codebox}[bash]
\begin{code}
brew install git
\end{code}
\end{codebox}

どちらの場合も、\cmd{git --version} でバージョンが表示されれば完了です。

\subsection{Linux}
\label{sec:install-linux}

お使いのディストリビューションのパッケージマネージャを使います。自分の環境に合うものを1つ選んで実行します。

\begin{codebox}[bash]
\begin{code}
# Debian / Ubuntu 系
sudo apt update && sudo apt install git

# Fedora / RHEL 系
sudo dnf install git

# Arch 系
sudo pacman -S git
\end{code}
\end{codebox}

\cmd{git --version} でバージョンが表示されれば完了です。

% ---------------------------------------------------------------------
\section{最初の設定}
\label{sec:first-config}

インストール直後に\textbf{一度だけ}行う設定です。コミットに記録される「あなたの名前とメールアドレス」を登録します。ここで設定した名前・メールが、後でGitHub上の履歴にも表示されます。

ターミナル（Windowsは Git Bash）を開き、次を順に実行します。

\begin{codebox}[bash]
\begin{code}
# あなたの名前を登録（GitHubのアカウント名に合わせるのがおすすめ）
git config --global user.name "your-name"

# あなたのメールアドレスを登録（GitHubで使う予定のメール）
git config --global user.email "you@example.com"

# 設定内容を確認する
git config --global --list
\end{code}
\end{codebox}

\begin{warnbox}
名前やメールにスペースが含まれる場合は、上のように \textbf{ダブルクォート \cmd{"} で囲んでください}。
\end{warnbox}

\cmd{--global} は「このPC全体の共通設定」という意味です。一度やれば、以降すべてのリポジトリに適用されます。

\begin{hintbox}
おすすめの追加設定（新規リポジトリの初期ブランチ名を \cmd{main} に統一）
\begin{codebox}[bash]
\begin{code}
git config --global init.defaultBranch main
\end{code}
\end{codebox}
\end{hintbox}

